\subsection{Efficiency and opportunities of our framework}\label{ss:zinc-crypto-efficiency-discussion}
The techniques used in Zinc+ in the context of SNARKs are relatively different from standard techniques that revolve around finite field arithmetic and polynomials with coefficients in said fields. We consider the problem of understanding how our framework compares to standard techniques an open problem. We provide an informal discussion of our observations around this question.

In terms of efficiency, we view Zinc+ as a framework that bootstraps the vast majority of SNARK costs to the single task of committing to multilinear polynomials with coefficients in $\QQ[X]$, $\ZZ[X]$, $\FF_q[X]$, etc.---in particular, for the applications we discussed, to polynomials with coefficients in $\BitPolyset{w}$. Indeed, in many ways, $\BitPolyset{w}$ acts as $w$ copies of a single problem to be solved. But on the other hand, the random projection technique onto finite fields compresses these $w$ copies into a single one, by introducing a small soundness error. Putting together this compression saving and the smaller initial arithmetizations that our UCS constraints produce, we observe in our implementation that the costs once in the finite field are tiny compared to the commitment costs. 

The fundamental question is then: \emph{does the tradeoff of smaller arithmetizations, plus the compression savings of projection, outweigh the cost of committing to these polynomials?} We believe that our preliminary benchmarks answer the question positively. 

We remark that committing to polynomials with coefficients in $\BitPolyset{w}$, while seemingly a ``$w$-wise'' task (especially because we commit to those using interleaved codes), can be implemented so that it does not scale linearly in $w$. Indeed, we observe engineering opportunities here by, for example, representing elements of $\BitPolyset{w}$ as \texttt{u32} types and using SIMD optimizations to operate on those. The fact that the Zinc+ framework concentrates all the computational burden on these commitments allows one to focus on optimizing a single type of operation. 

\subparagraph{Zip+, proof sizes, and future work} Zip+ is a Ligero/Brakedown-like IOPP (in modern instantiation) for $\QQ[X]$ and $\FF_q[X]$-valued multilinear polynomials.  Ligero/Brakedown on finite fields are known for having large proof sizes. And indeed our benchmarks show significant proof sizes relative to prover and verifier time (which on the other hand are very small). Besides being based on large-proof IOPPs, Zip+ experiments show a further proof size increase due to the fact that encodings of vectors in $\BitPolyset{w}$ scale linearly on $w$ in bit-size.  Ongoing work of ours is focused on designing a WHIR, Ligerito, Blaze, FICS/FACS, etc.\ \cite{whir, Ligerito, blaze, FICS-FACS} --IOPPs known for their small proofs among others-- variant for our framework. 



%\subparagraph{An example of the power of compression:} A paradigmatic example of how projecting onto a random field can save costs can be found in lookup constraints. One may have a lookup constraint into a very large table over some ring, e.g.\ $\ZZ[X]$ or $\FF_q[X]$ (for example in some verifiable FHE applications \cite{?}). As shown in \cref{?} the lookup predicate can be projected onto a random finite field $O(\lambda)$-bit-sized quotient of $\ZZ[X]$ or $\FF_q[X]$. One does not need to perform any 
     
